\documentclass{article}
\begin{document}
\section*{Native chosen reindexing: claim and boundary}
For the split endpoint interface, the left-associated positive tower now has a specified lift over each base arrow. The right-append step restricts the last grade along the destination of the recursively chosen prefix arrow. Its final component is the identity of the reindexed grade. The payload step makes the corresponding construction on the actual payload presheaf. A factorization proof computes the payload correction in an arbitrary competing arrow and proves uniqueness; no base-change equivalence is used to select the source objects or lifts.
\par Unique factorization gives a functor between the strict fibres, including the action on all vertical arrows. Repackaging verticality between the total over a fibre and the fibre of the total is strictly invertible. Identity and composite choices compare by unique invertible vertical transformations. Their comparisons respect vertical and horizontal composition and give the ordinary associator and unitors on rebased pasted routes. These are coherence witnesses, not equalities of the chosen functors.
\par Independent chain evaluation is extended to arrows over arbitrary base maps in the total of the iterate. Full hom equivalences and actual surjectivity on objects justify passage to the strict fibres. The square for each base arrow relates the actual native restriction to the iterate's separately constructed Grothendieck pull, naturally in vertical arrows. On canonical objects its target evaluates the presheaf map itself. Passing from these individual squares to a bundled comparison, including all mixed base and payload coherence, is not claimed. No assertion of normal strict splitting, the empty span, the length coproduct, flattening, or recovery of strict algebras follows.
\end{document}
